\scanpage{459}
\mspage{532}

All this being laid down, we then have the formulas
(160), (161), (161 bis), and the ``loop'' formula (157).

We now ask ourselves the question
under what conditions knowledge of
$(\alpha,\beta,f)$ (for example) among the quantities (153),
allows us to recover the quantities (152) ---
i.e. a section of $\mathcal M_{\rho,\sigma}$ --- and what conditions
must be imposed on $(\alpha,\beta,f)$. It
would suffice \inserted{(for uniqueness)} to be able to construct $\underline u$, whence
$\underline U,\underline a,\underline b$ by the first three formulas
(154). \deleted{\unsure{If, on the other hand we have}} We must
therefore find $\underline a\in\Gamma T_\rho$ such that $\alpha\underline a$ is a section
of $T_0$, --- indeterminacy in $\underline a$ is by
$\Gamma T_0\cap T_\rho$, and $\underline b\in\Gamma N'_\sigma$ such that $\beta\underline b$ is a section
of $T_0$ --- indeterminacy in \deleted{$\unsure{\Gamma N'_\sigma\cap}$} $\Gamma T_0\cap N'_\sigma$.
The indeterminacy in $\underline u$ is right
multiplication by a section of $T_0\cap T_\rho\cap N'_\sigma$, so it
will suffice that we have
\[
\text{(164)}\qquad T_0\cap T_\rho\cap N'_\sigma=\{1\}.
\]
\margin{in fact, we will \unsure{even} have $T_0\cap T_\rho=\{1\}$ in the cases \unsure{considered} in \uline{II}, \uline{III} and \unsure{following|sequel}\ldots}
Assuming this relation satisfied, what
conditions to put on $(\alpha,\beta,f)$ to come
from a \unsure{system} (152) --- i.e. from a section of $\mathcal M_{\rho,\sigma}$ ?

\mspage{458}

\scanpage{460}

We see that \unsure{if} a section \deleted{$\unsure{\varphi}$} $\varphi_0$ of $\mathcal M_{\rho,\sigma}$, \deleted{is}
\deleted{\unsure{replaced}} corresponding to $\alpha_0,\beta_0,f_0$, is replaced
by $\varphi_1\overset{\unsure{\text{definition}}}{=}f^{-1}.\varphi_0$, corresponding to $\alpha_1,\beta_1,f_1$, then
\[
\text{(165)}\qquad \alpha_1=\alpha_0,\quad\beta_1=\beta_0,\quad f_1=f_0f,
\]
for we see that if $\varphi_0$ corresponds to $(\underline U_0,\underline a_0,\underline b_0)$,
then $\varphi_1$ to $(U_1,\underline a_1,\underline b_1)$, then we have
\[
\text{(166)}\qquad U_1=f^{-1}U_0,\quad\underline a_1=\underline a_0,\quad\underline b_1=\underline b_0
\]
(for $\Theta_G\mid SG=$ the canonical inclusion $SG\hookrightarrow G$,
and $\Theta_\rho$\inserted{$,\Theta_\sigma$} \deleted{is} trivial on $SG$). \deleted{Thus This
shows that} \inserted{Thus \deleted{\unsure{it}}} we will have
\[
u_1=u_0\qquad (\text{for }\det U_1=\det U_0)
\]
whence (165). This shows that the conditions
sought on $(\alpha,\beta,f)$ depend only on
$\alpha,\beta$, not on $f$. We would like them
to be expressed by the validity of the loop
formula\inserted{(157)}, for a parameter \inserted{$\varepsilon,$} $\mu$ not specified $(\varepsilon^2=1)$
(which \unsure{will be} then uniquely determined --- and \unsure{will allow}
the \unsure{secondary definition} of \unsure{local notion|local action}\ldots).

\margin{to tell the truth, this condition is not \uline{sufficient} even in the cases considered --- we will find \unsure{therefore|in} a $\underline u\in B'_0$, and we must require that $\underline u\in T_0$ (already for $(\alpha,\beta)$ to be ``normalized'') i.e. $B=0$, where $B=ac+bd+bc$, if $\underline u=\left(\begin{smallmatrix}a&b\\c&d\end{smallmatrix}\right)$ (modular case) cf. proposition page~412'\ldots\ \unsure{To have it} $\unsure{R}$ ; \unsure{itself being} \unsure{additional condition}, \unsure{it|which} \unsure{is necessary to replace|is necessary to consider} $\mathcal M_{\rho,\sigma}$ \unsure{by the product} $\tfrac12$ \unsure{direct by} $L_0$, \unsure{on which we set} $\unsure{f_\Delta}=\unsure{0}$, \unsure{i.e.} $S\mathcal M_{\rho,\sigma}$.}

The sufficiency of this condition \unsure{will be}
the subject of the reflections in (II) (III),
\unsure{under} conditions (more draconian than those
\unsure{of the} preceding proposition (page~531') which we
have imposed on the triple $(\varepsilon_0,\rho,\sigma)$. It

\mspage{459}

\scanpage{461}
\mspage{533}

is possible that the conditions at which we
\inserted{\unsure{there}} arrived were more restrictive
than necessary, and that I shall be obliged
to come back to this later, in a context
more general than that of forms of
$\mathrm{Gl}(2)$, $\mathrm{GP}(1)$, $\mathrm{Sl}(2)$, but for the moment
the detailed study made by taking our inspiration
from the hypotheses brought out in \deleted{\unsure{the}} II, III will
have served to clear the ground for us, to clarify
our ideas on an important particular
case.

\uline{Remark.} Let us recall that we have seen that in
fact $\mu$ is already defined in terms of $\alpha$ alone,
and even in terms of $\alpha\rho(\alpha)^{-1}=\xi$ \unsure{;} \deleted{\unsure{indeed|therefore}}
But it is determined in terms of $\beta$
--- in fact, of $\beta\sigma(\beta)^{-1}$ --- only modulo
knowledge of $\varepsilon$ (in practice, modulo
sign\ldots). But if we impose (for example) $\varepsilon=1$,
the datum of $\alpha$ determines $\beta$ only
modulo \deleted{$\unsure{\beta}$} $\beta\longrightarrow\beta T_\sigma$\ldots

\mspage{460}

\scanpage{462}

The structure of $\mathcal M_{\rho,\sigma}$ being well understood,
let us return to the group $S\mathcal M_{\rho,\sigma}$\inserted{(cf. page~495')}, considered as
the semidirect product of $\mathcal M_{\rho,\sigma}$ by the subgroup
$L_0$ of $SG$, on which $\mathcal M_{\rho,\sigma}$ in fact operates
via $\chi$. Denoting by

\noindent (165)\quad
\diagram{\begin{tikzcd}[ampersand replacement=\&,column sep=large]
\deleted{\unsure{\varepsilon_0}:\unsure{U}}\;\inserted{\unsure{\mathcal E|\varepsilon}_{L_0}:G_a} \arrow[r,"\sim"] \&
\operatorname{Ker}\bigl(S\mathcal M_{\rho,\sigma}\longrightarrow\mathcal M_{\rho,\sigma}\bigr)
\end{tikzcd}}

\inserted{corresponding to $\varepsilon_0$,}
we will have by \deleted{\unsure{definition}} construction
\[
\text{(166)}\qquad
U\bigl(\unsure{\mathcal E|\varepsilon}_{L_0}(\lambda)\bigr)
=\unsure{\mathcal E|\varepsilon}_{L_0}(\mu\lambda)
\qquad\text{where }\mu=\chi(U).
\]

\deleted{The group $S$}

\mspage{461}

\scanpage{463}
\mspage{534}

\deleted{\unsure{XI}} \textcircled{$9^\circ$} \uline{The subgroups}
$\mathcal M_{\rho,\sigma}(0)$,
\deleted{$\mathcal M_{\rho,\sigma}(\tfrac12),\mathcal M_{\rho,\sigma}(-\unsure{\overline J|\overline\jmath})$.}
\inserted{$\mathcal M_{\rho,\sigma}[0]$, and the group}
$S\mathcal M_{\rho,\sigma}$, the isomorphism
\diagram{\begin{tikzcd}[ampersand replacement=\&,column sep=large]
S\mathcal M_{\rho,\sigma} \arrow[r,"\sim"] \& \Sigma^+_{\rho,\sigma}
\end{tikzcd}}

The group $\mathcal M_{\rho,\sigma}$ is an extension of $\mathcal T_{\rho,\sigma}$ by
$\unsure{\widetilde G|\widetilde{\mathcal C}}^+_{\rho,\sigma}\simeq SG$, and the three subgroups which we are going
to make explicit will be \deleted{sections}\inserted{\unsure{subgroups|subgroups}}, \deleted{partial} with the
grain of salt that for the third, it will be a
partial section only, over the
subgroup \inserted{$\mathcal T'_{\rho,\sigma}$} of index $2$ in $\mathcal T_{\rho,\sigma}$, kernel of
$\varepsilon_3$.

$[\![$Let us determine in a general way the sections
of $\mathcal M_{\rho,\sigma}$ over $\mathcal T_{\rho,\sigma}$. Since we have

\noindent (1)\quad
\diagram{\begin{tikzcd}[ampersand replacement=\&,column sep=large,row sep=small]
R_{\rho,\sigma} \arrow[r] \arrow[d,phantom,"{\unsure{\wr|\simeq}}" description] \& \mathcal T_{\rho,\sigma} \quad \text{epimorphism with kernel }\mu_2\\
S'T_\rho\times SN'_\sigma\times G_m \&
\end{tikzcd}}

\margin{in the end this won't be of much use\ldots}

it comes to the same to give oneself a section,
or: a homomorphism

\noindent (2)\quad
\diagram{\begin{tikzcd}[ampersand replacement=\&,column sep=large]
R_{\rho,\sigma} \arrow[r,"{\deleted{\unsure{\Xi|\zeta}}\inserted{\psi}}"] \& \mathcal M_{\rho,\sigma}
\end{tikzcd}}

which is \deleted{trivial}
\[
\text{(3)}\qquad
\left\{\begin{aligned}
\text{a)}\quad &\text{trivial on }\mu_2\\
\text{b)}\quad &\deleted{\unsure{\Xi|\Sigma}(x)\unsure{\equiv|=}x\ \bigl(\unsure{\widetilde G|\widetilde{\mathcal C}}^+_{\rho,\sigma}\bigr)\quad\text{identically}.}\quad\inserted{\text{``lifts'' (1)}}
\end{aligned}\right.
\]

The homomorphisms $\unsure{\Xi|\Sigma}$ satisfying (3) correspond
to the splittings of the \uline{trivial} extension of
$R_{\rho,\sigma}$ by $SG$ (\uline{product} extension), hence
simply to the \uline{homomorphisms}

\noindent (4)\quad
\diagram{\begin{tikzcd}[ampersand replacement=\&,column sep=large]
R_{\rho,\sigma} \arrow[r,"{\deleted{\unsure{\Xi|\zeta}_{\unsure{0}}}\inserted{\psi^0}}"] \& SG
\end{tikzcd}},

\mspage{462}

\scanpage{464}

by associating to $\psi^0\deleted{\unsure{\xi_0}}$ the homomorphism
\[
(5)\qquad x\longmapsto \deleted{\unsure{\xi}}\psi(x)
=x\deleted{\unsure{\xi}}\psi^0(x)
=\deleted{\unsure{\xi(x)x}}\psi^0(x).x
\]
Thus condition a takes the form
\[
(6)\qquad \deleted{\unsure{\xi}}\psi^0_{\unsure{0}}(x)=x
\quad\text{on }\mu_2=R_{\rho,\sigma}\cap SG.
\]

\margin{\deleted{\unsure{$\rho$}}}
The $\psi,\psi^0$ corresponding to the \deleted{\unsure{``three''}} sections
$\mathcal M_{\rho,\sigma}^{(0)}$, $\mathcal M_{\rho,\sigma}^{(\frac12)}$,
$\mathcal M_{\rho,\sigma}^{(-\overline J)}$ will be denoted resp.\textsuperscript{ly}
\[
(\psi_{\unsure{0}},\psi^0_0),\ (\psi_{1/2},\psi^0_{1/2}),\
(\psi_{-\overline J},\psi^0_{-\overline J}).\qquad\big]\!\big]
\]

1°) \uline{Case of} $\mathcal M_{\rho,\sigma}^{(0)}$. With the notations
of page 529, it \unsure{will correspond} \deleted{\unsure{$\underline U$ and the}}
\inserted{\unsure{$\underline u$}} of $\mathcal M_{\rho,\sigma}$ satisfying one of the
\uline{equivalent} conditions
\[
(7)\qquad f=1,\quad\inserted{\underline U=\underline u,}\quad
\underline U\in\Gamma T_0,\quad \alpha_0=\alpha,\quad\beta_0=\beta
\]
so we have
\[
(8)\qquad \mathcal M_{\rho,\sigma}^{(0)}=\Theta^{-1}(T_0).
\]
\inserted{(in terms of the description of $\mathcal M_{\rho,\sigma}$ on page 526' and following)}

Let us check that this is a \unsure{subgroup-section}.
We have
\[
\overbrace{\mathcal M_{\rho,\sigma}^{(0)}\cap\widetilde{\unsure{G|\mathcal C}}_{\rho,\sigma}^{+}}
^{\inserted{(\Theta|SG)^{-1}(T_0)}}=\{1\}
\]
\margin{\unsure{and}}
for $\Theta|(\widetilde{\unsure{G|\mathcal C}}_{\rho,\sigma}^{+}=SG)$ is the canonical
inclusion $SG\hookrightarrow G$, now $SG\cap T_0=\{1\}$.
It remains to check that
\diagram{\begin{tikzcd}[ampersand replacement=\&]
\mathcal M_{\rho,\sigma}^{(0)} \arrow[r] \& \mathcal T_{\rho,\sigma}
\end{tikzcd}}
(which is a monomorphism by \unsure{$(a)$|$(\alpha)$}) is an epimorphism. But in

\mspage{463}

\scanpage{465}
\mspage{535}

terms of the canonical \unsure{realization}
\[
(9)\qquad\mathcal M_{\rho,\sigma}\simeq
\underbrace{R_{\rho,\sigma}\times SG}_{\widetilde{\mathcal M}_{\rho,\sigma}}
/\underset{\text{diagonal}}{\mu_2}
\]
and of the homomorphism
\margin{\diagram{\begin{tikzcd}[ampersand replacement=\&]
\kappa:R_{\rho,\sigma}\arrow[r]\&\mathbf G_m
\end{tikzcd}}
the canonical homomorphism \unsure{making it possible} to express $\Theta|R_{\rho,\sigma}$}
\diagram{\begin{tikzcd}[ampersand replacement=\&]
(10)\qquad i_{T_0}:\mathbf G_m\arrow[r,"\sim"]\&T_0\arrow[r,hook]\&G
\end{tikzcd}}
\unsure{isomorphism|homomorphism} inverse to the isomorphism
\diagram{\begin{tikzcd}[ampersand replacement=\&]
\delta_{T_0}=\det|T_0:\quad T_0\arrow[r,"\sim"]\&\mathbf G_m,
\end{tikzcd}}
we have, for an element
\diagram{\begin{tikzcd}[ampersand replacement=\&,column sep=1em,row sep=.2em]
(\underline r,\&U_0)\\
\inserted{R_{\rho,\sigma}}\arrow[u]\&\inserted{SG}\arrow[u]
\end{tikzcd}}
of $\Gamma\widetilde{\mathcal M}_{\rho,\sigma}$
\[
\Theta(\underline r,U_0)=\kappa(\underline r)U_0,\quad
\text{determinant }\kappa(r)^2
\]
to say that this is $\unsure{\in U}$ therefore means we have
\[
\kappa(\underline r)U_0=i_{T_0}(\kappa(\underline r)^2)
\quad\text{i.e. }\deleted{\unsure{U_0=}}
\]
\[
(11)\qquad U_0=\kappa(\underline r)^{-1}i_{T_0}(\kappa(\underline r)^2)
\]
so the inverse image $\widetilde{\mathcal M}_{\rho,\sigma}^{(0)}$ of
$\mathcal M_{\rho,\sigma}^{(0)}$ in $\widetilde{\mathcal M}_{\rho,\sigma}$ consists of the
\[
(12)\qquad
\bigl(\underline r,\kappa(\underline r)^{-1}i_{T_0}(\kappa(\underline r))^2\bigr)
\in\Gamma\widetilde{\mathcal M}_{\rho,\sigma}
\]
in other words this corresponds to a \unsure{group} section, with
\[
(13)\qquad\psi^0_0(\underline r)=\kappa(\underline r)^{-1}i_{T_0}(\kappa(\underline r)^2).
\]
As it should be, we do indeed have here a \uline{homomorphism}
\diagram{\begin{tikzcd}[ampersand replacement=\&]
\psi^0_0:R_{\rho,\sigma}\arrow[r]\&SG
\end{tikzcd}}
and its value on $\mu_2$ is the identity,
for $\kappa(\underline r)^2=1$ and we have
\[
\psi^0_0(r)=\kappa(\underline r)^{-1}=
\begin{gathered}\kappa(\underline r)\\\Vert\\\underline r\end{gathered}
\]

\uline{Remark}\quad\deleted{Before \unsure{qui}}
To define this section, it sufficed

\mspage{464}

\scanpage{466}

(from the point of view of pure calculation) to give oneself an arbitrary homomorphism, denoted
$\unsure{i'_{T_0}|i_{T_0}}$, from $\mathbf G_m$ to $SG$, and to set (13).

Before leaving $\mathcal M_{\rho,\sigma}^{(0)}$, I am going to introduce the
slightly larger group (of dimension 4, instead of 3)
\[
(14)\qquad\mathcal M_{\rho,\sigma}^{[0]}
\overset{\mathrm{dfn}}{=}\Theta^{-1}(B'_0)
\]
where, let us recall, $B'_0$ is defined as the $\frac12$ direct product
\[
(15)\qquad B'_0=T_0.L_0
\]
This time, $U=U_0\kappa(r)\inserted{\in\Gamma B'_0}$ is not determined
by knowledge alone of its determinant\inserted{\unsure{$2$}}
$\underline\kappa(r)^2$, which determines it only modulo $L_0$.
This corresponds to the fact that
$\underset{\inserted{=B'_0\cap SG}}{
\mathcal M_{\rho,\sigma}^{[0]}\cap
\overbrace{\widetilde{\unsure{G|\mathcal C}}_{\rho,\sigma}^{+}}^{\inserted{SG}}}$
is no longer reduced to $(1)$, but to $L_0$, which therefore appears as
an invariant subgroup of $\mathcal M_{\rho,\sigma}^{[0]}$, and we obviously \unsure{will have|have}
(taking into account that $\mathcal M_{\rho,\sigma}^{(0)}\subset\mathcal M_{\rho,\sigma}^{[0]}$ and
\diagram{\begin{tikzcd}[ampersand replacement=\&]
\mathcal M_{\rho,\sigma}^{(0)}\arrow[r,"\sim"]\&\mathcal T_{\rho,\sigma}
\end{tikzcd}}) \unsure{by|that}
\[
(16)\qquad\begin{aligned}
\mathcal M_{\rho,\sigma}^{[0]}&=\mathcal M_{\rho,\sigma}^{(0)}.L_0\\
&\qquad\text{product }\frac12\text{ direct}
\end{aligned}
\]
This is again a construction which makes sense as soon as we have a $T_0$ as above,
\deleted{\illegible{2}} and a subgroup \deleted{\uline{\unsure{affine}}} $L_0$ of $SG$,

\mspage{465}